\subsection{A uniform cut criterion in any degree}
\label{sec:uniformcutdecay}

Intersecting critical curves need not admit a simple formal normal
form. The following criterion avoids choosing one. Its hypothesis
concerns the full phase $tP-q$, uniformly in large complex $t$;
bounded critical points alone would not suffice.

\begin{theorem}\label{thm:uniformcuts}
Let $P,q\in\C[x_1,\ldots,x_n]$, with $P$ nonconstant. Suppose
there are fixed $c,R_0,t_0>0$ such that
\begin{equation}\label{eq:uniformphasegradient}
 \|\nabla_x(tP-q)(x)\|\ge c\|x\|
       \quad(|t|\ge t_0,\ \|x\|\ge R_0).
\end{equation}
Along every fixed ray $t=Re^{i\theta}$, the full rapid-decay
homology has a flat cut basis $\Gamma_i(t)$, attached to the
critical values $c_i(t)$ of $tP-q$, such that for every fixed
polynomial $a$,
\begin{equation}\label{eq:normalizedcutdecay}
 e^{-c_i(t)}\int_{\Gamma_i(t)}a(x)e^{tP-q}\,
                 dx_1\cdots dx_n\longrightarrow0.
\end{equation}
Each $c_i(t)=\alpha_i t+O(1)$, where $\alpha_i$ is a critical
value of $P$. If all critical values of $P$ lie in a closed
half-plane through zero, a supporting $t$ ray makes every flat
period of every such fixed form tend to zero.
\end{theorem}

We give the completeness, tail and limiting-form arguments separately.
None follows just from counting the phase-critical points.

\begin{proof}
For every sufficiently large $t$, \eqref{eq:uniformphasegradient}
makes the phase-critical set finite and contained in one fixed
ball. It also gives complex tameness at infinity. Pham's tame
Milnor decomposition~\cite[Part~I and Appendix]{Pham} identifies
the complete middle rapid-decay homology with the local vanishing
homologies carried along cuts from its critical values. The rank
is the sum of their Milnor numbers, including degenerate points.
If this rank is zero the conclusion is immediate. Otherwise
retain all these local generators; identical critical values
are grouped, not discarded.

On the chosen ray, the critical points and values are algebraic
functions of $R$, after a finite cover. Enlarge $R_1$ so branches
and multiplicities stabilize. Choose $|\eta|=1$ with
$\Re\eta\ge1/2$ so the cuts
\[
 S_t(x)=c_i(t)-\eta s,\qquad s\ge0,\quad S_t=tP-q,
\]
are disjoint for $R>R_1$, except at identical grouped endpoints.
There are only finitely many forbidden choices: every nonzero
difference of two algebraic critical values has one limiting
direction. The cut angle is independent of $\theta$. In
particular a supporting $t$ ray need not be a generic ray.

Along a cut, normalized complex gradient transport is
\[
 x_s=-\eta\,\overline{\nabla S_t}/\|\nabla S_t\|^2,
               \qquad S_t(x(s))_s=-\eta.
\]
Outside the fixed ball,
$|d\|x\|^2/ds|\le2/c$. All local vanishing generators start
in that ball. With fixed sufficiently large $A$ and $B>2/c$,
they therefore have representatives inside the carrier
\begin{equation}\label{eq:uniformcarrier}
 \|x\|^2\le A+Bs.
\end{equation}
At its outer boundary the transport points strictly inward
relative to this inequality. No trajectory escapes at finite
cut length, and the cuts avoid all other critical values.

We must control a single family of representatives. The gradient
flow itself need not be semialgebraic. Instead compactify $x$
by $x/\sqrt{1+\|x\|^2}$ and $s$ by $s/(1+s)$. Take the closure
of each constrained cut carrier and mark its rapid boundary
$G$, where $s=\infty$. Inequality~\eqref{eq:uniformcarrier}
forbids spatial infinity at finite $s$. For fixed $R$, the
relative chains of this compact pair are precisely locally
finite cut chains with their ends allowed only at rapid
infinity. This follows also by a compatible triangulation of
the compact pair: its relative simplices restrict to the usual
locally finite chain model in the open carrier. The exponential
there is exactly $e^{c_i-\eta s}$.

For the full unconstrained cut carrier, mark and remove any
spatial boundary at finite $s$; it is not a rapid boundary.
The constrained carrier, the full carrier, their rapid boundaries
and this excluded boundary form a finite family of semialgebraic
subsets in the same compactification. Apply the compatible
finite-family form of Hardt triviality to projection onto $R$.
It gives one trivialization on an unbounded interval, preserving
these subsets and their inclusions. The precise form is stated
in Drusvyatskiy--Lewis~\cite[Theorem~2.25]{DrusvyatskiyLewis};
it is Hardt's theorem. Semialgebraic triangulation and the
chain-current boundary rule are as in
Hardt--Lambrechts--Turchin--Voli\'c~\cite[Sections~2.1 and~3,
especially Theorem~3.5]{HLTV}.

Gradient transport already showed that the constrained inclusion
contains every local vanishing generator. At one parameter choose
finitely many integral relative chains whose images are a basis
over $\C$ of the full rapid-decay space. A nonprimitive integral
image causes no problem over $\C$. Transport these chains by
the compatible trivialization. They form a jointly semialgebraic
family with fixed integer multiplicities and supports satisfying
\eqref{eq:uniformcarrier}. Their images are a flat basis of the
actual homology local system: on each compact parameter interval
the transport preserves the relative pairs and is proper on
every finite phase band. Thus this is a completeness statement
for actual cycles, not an identification of presumed formal blocks.

Here is the needed uniform mass bound. In a fixed semialgebraic
family of $d$-chains with bounded multiplicities, the mass inside
a ball of radius $L\ge1$ is at most $C L^d$, uniformly in the
parameters. Partition tangent planes by their largest coordinate
$d$-minor. The Cauchy--Binet identity gives a lower bound
$\binom Nd^{-1/2}$ for some such minor in ambient dimension $N$.
The area formula then bounds that piece by its coordinate
projection volume times the number of preimages. Semialgebraic
uniform finiteness bounds those finite fibers independently of
the parameters; see, for example,
\cite[Theorem~2.32]{DrusvyatskiyLewis}. Sum the finitely many
pieces and multiplicities. Apply the same argument to boundary
chains. This proves uniformity, rather than merely pointwise
finite volume.

Lift the cut chains to $(x,s)$. Their bands $j\le s\le j+1$
lie in balls of radius $O(1+j)$ by
\eqref{eq:uniformcarrier}. Their fixed semialgebraic descriptions
and the mass estimate give, with constants independent of $R$,
\begin{equation}\label{eq:uniformcuttail}
 \left|e^{-c_i(t)}\int_{s\ge L}a e^{S_t}\,dx_1\cdots dx_n\right|
 \le C\sum_{j\ge\lfloor L\rfloor}(1+j)^M e^{-j/2}.
\end{equation}
The polynomial amplitude changes only $M$. This tail tends to
zero uniformly in $R$.

It remains to bound a fixed phase band. A bounded critical branch
$x(r)$, $r=1/t$, satisfies
\[
 q(x(r))=-c_i(r)-r c_i'(r).
\]
At a stationary point the chain rule gives $c_i'(r)=-P(x(r))/r^2$;
eliminating $P$ gives this identity. Since $q(x(r))$ is
bounded, its Puiseux expansion forces the only possible polar
term of $c_i$ to be $\alpha_i/r$. Its coefficient is
$\alpha_i=P(x(0))$, and $\nabla P(x(0))=0$. Hence
$c_i(t)=\alpha_i t+O(1)$ even at critical collisions.

Choose a generic fixed truncation level $L$. The truncated lifted
chains have compact supports and uniformly bounded mass and
boundary mass. The latter follows from the same semialgebraic
estimate applied to their slices; after a finite stratification
there is an arbitrarily large sequence of such levels. By the
Federer--Fleming compactness theorem, every sequence $R\to\infty$
has a subsequence converging to an integral $n$-current. We use
the version with both mass and boundary-mass bounds,
Simon~\cite[Chapter~6, Section~3, Theorem~3.15]{SimonGMT}.
Its rectifiability conclusion is essential.

On the lifted chain,
\[
 P(x)-\alpha_i=
       \frac{c_i(t)-\alpha_i t+q(x)-\eta s}{t}.
\]
Every limiting current is therefore supported in
$\{P=\alpha_i\}\times[0,L]$. Stratify this complex algebraic
level set into smooth complex strata. Almost everywhere for the
rectifiable limiting current, its tangent plane projects into
the tangent of one of these strata. Each has complex dimension
less than $n$. The holomorphic $n$-form $dx_1\cdots dx_n$ thus
vanishes on that tangent plane. Lower-dimensional strata carry
no $n$-mass; the same argument covers singular strata of larger
real dimension. Testing with the fixed smooth form
$e^{-\eta s}a(x)dx_1\cdots dx_n$, with a compact cutoff, gives
zero on every subsequential limit. Consequently the normalized
integral on $0\le s\le L$ tends to zero. This does not assert
that the real volume of the chains tends to zero.

First let $R\to\infty$, then let $L\to\infty$ through the
allowed levels in \eqref{eq:uniformcuttail}. This proves
\eqref{eq:normalizedcutdecay}. Finally choose $\theta$ with
$\Re(e^{i\theta}\alpha)\le0$ for all critical values $\alpha$
of $P$. The factors $e^{c_i(t)}$ are then bounded, so every basis
period tends to zero. A flat cycle is a constant linear
combination of this full basis, proving the last assertion.
\end{proof}

\begin{corollary}\label{cor:twocriticalvalues}
Let $q$ be a full affine Gaussian quadratic and assume
\eqref{eq:uniformphasegradient}. If $P$ has at most two critical
values, every original marked Gaussian period is either zero or
is not a finite exponential-polynomial. In particular $P$ is
neither pure nor safe.
\end{corollary}
\begin{proof}
Any two complex numbers lie in a closed half-plane through zero.
For two nonzero values normalize them to $z_1,z_2$ on the unit
circle. If $z_1+z_2\ne0$, use the direction
$-\overline{z_1+z_2}/|z_1+z_2|$; if their sum is zero, use
$i\overline{z_1}$. Zero or missing values add no condition.

A nonzero constant original marked period would, by the rational
seed restriction and actual-period extension in
Lemma~\ref{lem:seedfunctional}, be represented by a flat period
of that fixed polynomial form. Its generic exponential system
has $f=t$, $g=tP-q$ on smooth affine space and regular input
$\mathcal O$; $f,g$ are algebraically independent because $P$
is nonconstant. Remove the finite exceptional parameter set of
the Hien--Roucairol comparison. Theorem~\ref{thm:uniformcuts}
makes the resulting constant period tend to zero, a contradiction.
No convergence of the formal moment series was assumed.

A constant-coefficient operator reduces a nonzero finite
exponential-polynomial to $c e^{\beta t}$ while replacing its
original marker by another fixed polynomial in $P$ times that
marker. Shift $P$ by $-\beta$ and normalize. Its gradient is
unchanged and its critical values are still at most two. The
same contradiction applies. Since $\Phi_1(0)=1$, it cannot be
a polynomial, proving also the safety exclusion.
\end{proof}
